\subsection{Reverse excision and the three meanings of fibration}

\begin{claim}{B.6}{Refuted}{The reverse Mayer--Vietoris axioms}
For a derivator $\D$ with $\D(e)\not\simeq e$, the class
$\Wlc$ satisfies neither Loc~(6b) nor Loc~(6c) of Chapter XVI.
If $\D(e)\simeq e$, all arrows belong to $\Wlc$ and both
axioms hold. For $\Wone$ of a nontrivial derivator,
Loc~(6b) always fails, whereas Loc~(6c) holds exactly when
$\D(e)$ is a preorder up to equivalence.
Thus the unqualified Loc~(6), which includes
Loc~(6b), fails for every nontrivial coefficient theory, even
for spaces and ordinary weak homotopy equivalences. In contrast,
each constant homology localizer $W_E$ satisfies all three
Mayer--Vietoris implications Loc~(6a,b,c).
\end{claim}
\begin{proof}
Write a Mayer--Vietoris square in the order
$(X_0,X_1,X_2,X_3)$, where $X_1,X_2$ are full sieves of
$X_3$, their union is $X_3$, and $X_0=X_1\cap X_2$.
For a map of two such squares, the manuscript's precise
implications are
\[
\begin{array}{ll}
 (6b):& f_1,f_2,f_3\in W\Longrightarrow f_0\in W,\\
 (6c):& f_0,f_3\in W\Longrightarrow f_1,f_2\in W.
\end{array}
\tag{B.6.1}
\]
We give separate countercubes, so no implication from (6c)
to (6b) is being assumed.

For any small category $A$, let $C_v=A\star\{v\}$ and
$C_w=A\star\{w\}$ be its two terminal cones and put
$\Sigma A=A\star\{v,w\}$, with no arrows between $v$ and
$w$. Its objects are those of $A$ and $v,w$. Its arrows
are those of $A$, the two new identities, and exactly one
arrow $a\to v$ and one arrow $a\to w$ for every $a\in A$.
Composing any arrow of $A$ with either terminal arrow gives
the terminal arrow from its source; these and the compositions
in $A$ specify the entire composition law. The inclusions
$C_v,C_w\subset\Sigma A$ are full sieves with intersection
$A$. Their nerves form a homotopy pushout, and
$B\Sigma A\simeq\Sigma BA$ is the unreduced suspension.
It suffices to find
\[
 A\longrightarrow e\notin\Wlc,
 \qquad \Sigma A\longrightarrow e\in\Wlc.
\tag{B.6.2}
\]

First suppose that $\D(e)$ has distinct parallel morphisms
$\alpha,\beta:M\to N$. Choose a nontrivial integrally acyclic
group $G$ as constructed in B.10, and take $A=BG$, the
one-object category with arrows $G$ and group multiplication
as composition. The CW space $BG$ is connected and acyclic.
Its suspension is simply connected and acyclic, and therefore
weakly contractible by homological Whitehead. Thus the second
condition of (B.6.2) holds by A.2.

For the first condition, let $a:e\to BG$ be its object.
The coherent object $a_!M\in\D(BG)$ has value
$\coprod_{g\in G}M$ at that object; its $G$-action is regular
permutation of the summands. This is a pointwise Kan-extension
calculation, and constructs a coherent object rather than
merely postulating a diagram in $\D(e)$. A nonidentity
permutation is detected by a map to $N$ assigning $\alpha$
to one summand and $\beta$ to all the others. Consequently
its action is nontrivial. It cannot be isomorphic to a
constant object, whose action is trivial. Every object of
$\D(BG)$ is locally constant since every arrow of $BG$ is
invertible. Thus $\D(e)\to\LC_\D(BG)$ is not essentially
surjective, proving the first condition.

It remains to treat a nontrivial $\D(e)$ which is a preorder
up to equivalence. For every nonempty group $G$, a nonempty
coproduct of copies of an object $M$ of this category is
isomorphic to $M$. The counit
$a_!a^*F\to F$, for $F\in\D(BG)$, is therefore an
isomorphism on its sole point evaluation, and hence an
isomorphism by Der~2. Likewise the canonical map
$a_!M\to p^*M$, where $p:BG\to e$, is an isomorphism.
The adjunction for $a$ now gives
\[
 \Hom_{\D(BG)}(p^*M,p^*N)
 \simeq\Hom_{\D(BG)}(a_!M,p^*N)
 \simeq\Hom_{\D(e)}(M,N).
\]
Thus every $BG\to e$ belongs to $\Wlc$. Take instead
$A=e\amalg e$. Its projection is not in $\Wlc$: the
diagonal $\D(e)\to\D(e)^2$ is not essentially surjective
because $\D(e)$ has two nonisomorphic objects. The category
$\Sigma A$ is the four-object poset
$a,b<v,w$. Its nerve is a circle, hence is weakly equivalent
to $B\Z$, and the preceding argument proves the second
condition of (B.6.2). Explicitly, a weak equivalence
$\Sigma A\to B\Z$ sends $a\to v$ to $1\in\Z$, the
other three nonidentity arrows to $0$, and every object to
the unique object; its map on the unique cycle has degree one.

For Loc~(6b), map the square
$(A,C_v,C_w,\Sigma A)$ to the identical double-cone
construction on $e$. Each cone has a terminal object, and
both total categories are $\Wlc$-aspherical by (B.6.2).
Thus $f_1,f_2,f_3$ belong to $\Wlc$, whereas
$f_0:A\to e$ does not. For Loc~(6c), use the map
\[
 (A,A,C_w,C_w)\longrightarrow(A,C_v,C_w,\Sigma A).
\tag{B.6.3}
\]
Both quadruples are Mayer--Vietoris squares, allowing one
sieve to equal the whole category. Here $f_0=\id_A$,
$f_2=\id_{C_w}$, and $f_3:C_w\to\Sigma A$ belong to
$\Wlc$, but $f_1:A\to C_v$ does not, by two-out-of-three
with $C_v\to e$. This directly disproves (6c).
The finite thin-coefficient instance, including all objects,
arrows, compositions, functoriality and sieve checks, is
verified by \texttt{code/B\_verify\_reverse\_excision.py};
its complete output is included.

For the asserted $1$-complete comparison, first assume
that $\D(e)$ is nonthin. Reverse the intersection map in
the first cube: map the double cone of $e$ to the double
cone of the acyclic $BG$. The three cone and total maps
are weak homotopy equivalences. But restriction along
$e\to BG$ is not full: with the objects $M,N$ chosen
above, it induces the nonbijective diagonal
\[
 \Hom(M,N)\longrightarrow\Hom(M,N)^G
\]
on morphisms from the free coherent $G$-object $a_!M$
to the constant object $N$. This refutes (6b).

For (6c), choose a nontrivial acyclic group $H$ and,
by B.10, an acyclic simple group $G$ containing $H$
properly. Its larger cardinality guarantees properness
of the inclusion. Let $Y_1$ be the categorical mapping
cylinder of $BH\to BG$. It has two objects $h,g$,
endomorphism groups $H,G$, and
$\Hom(h,g)=G$, with $\Hom(g,h)=\varnothing$.
Composition on the cross-arrows is left multiplication
by $G$ and right multiplication by the specified image
of $H$. This specifies all its arrows and compositions.
The full subcategory $BH$ is a sieve; projection
$Y_1\to BG$ is a weak homotopy equivalence, as seen
from its terminal-index Grothendieck construction.
Let $C_H=BH\star\{w\}$ and glue it to $Y_1$ along
$BH$, giving $Y_3$. Besides the arrows just listed,
$Y_3$ has $\id_w$ and a unique $h\to w$, fixed by
precomposition with $H$, and no arrows $g\to w$ or
out of $w$. Its two pieces $Y_1,C_H$ are full sieves.
Its classifying space is the homotopy cofiber of
$BH\to BG$. Van Kampen gives
$\pi_1BY_3=G/\langle\!\langle H\rangle\!\rangle=1$
by simplicity; Mayer--Vietoris gives zero reduced
integral homology since both groups are acyclic.
Thus $BY_3$ is weakly contractible.

Use the map of squares
\[
 (BH,BH,C_H,C_H)\longrightarrow(BH,Y_1,C_H,Y_3).
\tag{B.6.5}
\]
Its $f_0$ is an identity and its $f_3$ is a weak
homotopy equivalence. Its $f_1$ is not $1$-complete.
Indeed, restriction of the free coherent $G$-object
to $BH$ is a coproduct of free coherent $H$-objects,
one for each coset in $G/H$. To see this at the
coherent level, choose coset representatives, use
the corresponding arrows of $BG$ to construct the
maps from the free $H$-objects by Kan adjunction,
and apply point evaluation: their disjoint regular
orbits partition $G$, so that map is an isomorphism
by Der~2. The morphism comparison to the constant
object $N$ is consequently
\[
 \Hom(M,N)\longrightarrow\Hom(M,N)^{G/H},
\]
again a nonbijective diagonal. Since $Y_1\to BG$
is a weak equivalence, the same failure holds for
$BH\to Y_1$. This proves failure of (6c) for
$\Wone$ with nonthin coefficients, including all
the nonzero stable coefficient categories.

There is a different answer for thin coefficients.
If $P=\D(e)$ is a nontrivial preorder up to equivalence,
then every $\D(A)$ is thin as well: its diagonal
$F\to F\times F$ is pointwise an isomorphism,
hence an isomorphism by Der~2, which forces every
Hom set to have at most one element. Let
$u:\operatorname{Ob}(A)_{\mathrm{disc}}\to A$ be
the inclusion of its discrete object set. The
unit $F\to u_*u^*F$ is a pointwise isomorphism.
At $a$ the right-hand side is the product of the
values $F_b$ over arrows $b\to a$, in presheaf
variance. The value $F_a$ maps to every factor by
the underlying diagram, and occurs as the identity
factor; in a preorder these two facts make the
unit an isomorphism. The Kan adjunction therefore
makes the underlying-diagram functor fully faithful.
Conversely, for a presheaf of values in $P$, applying
$u_*$ to its values has exactly those point values
by the same product calculation. Thus
$\D(A)\simeq\Fun(A^{\op},P)$, and
\[
 \LC_\D(A)\simeq P^{\pi_0A},\qquad
 \Wone=\{f:\pi_0f\text{ is surjective}\}.
\tag{B.6.6}
\]
Here the last equivalence uses nontriviality of the
complete preorder $P$: an omitted target component
is detected by its distinct initial and terminal
objects. Map the double cone of the empty category
to the double cone of $e$. Its intersection map
$\varnothing\to e$ is not in $\Wone$; both cone
maps are weak equivalences, and the total map from
two points to a connected category is surjective
on components. This refutes (6b).

To verify (6c) in this case, suppose $f_0,f_3$
are surjective on components and let $C$ be a
component of $Y_1$. If $C$ meets $Y_0$, surjectivity
of $\pi_0f_0$ gives an object of $X_0\subset X_1$
mapping into $C$. If it does not, the sieve-cover
description of the nerve shows that $C$ is a whole
component of $Y_3$ disjoint from $Y_2$.
Surjectivity of $\pi_0f_3$ then gives an object
of $X_3$ mapping into $C$; it must lie in $X_1$,
since an object of $X_2$ maps into $Y_2$.
Thus $\pi_0f_1$ is surjective, and the symmetric
argument proves the same for $f_2$. This proves
the exact thin-coefficient exception for (6c).

If $\D(e)\simeq e$, G.1 shows $\D(A)\simeq e$ for
every $A$, and the exception in the statement follows.
Finally, a sieve cover gives a homotopy-pushout square of
nerves. Applying $E\wedge\Sigma^\infty_+$ gives the exact
triangle
\[
 E\wedge\Sigma^\infty_+BX_0\longrightarrow
 (E\wedge\Sigma^\infty_+BX_1)\oplus
 (E\wedge\Sigma^\infty_+BX_2)\longrightarrow
 E\wedge\Sigma^\infty_+BX_3.
\tag{B.6.4}
\]
In a map of exact triangles, equivalence of any two terms
implies equivalence of the third. An equivalence of the
displayed block-diagonal middle map implies equivalence of
each block, since each is a retract. This proves (6a), (6b)
and (6c) for $W_E$, in particular for the constant class
of $D(k)$. This stable excision argument does not apply to
the $2$-complete class of $D(k)$.
\end{proof}

\begin{claim}{B.7}{Proved}{Classification of the fibration axioms}
Let $W$ be a fundamental localizer with
$W\subseteq W_0$. Define $\mathcal F_W$ to be the class
of Grothendieck fibrations whose transport functors belong
to $W$, and define $\operatorname{Fib}_W$ by preservation
of all $W$-equivalences under strict categorical pullback.
Then each of Loc~(7a), Loc~(7b) and Loc~(7c) holds if and
only if $W=W_\infty$. In particular this applies to
$W=\Wlc$ for every nontrivial derivator, by B.3 and G.1.
For the trivial derivator all three axioms hold.

If $W$ is accessible, Loc~(7F), with the manuscript's
class $\operatorname{Fib}_W$, holds if and only if $W$
is proper, equivalently if and only if it is generated
as a fundamental localizer by a set of projections
$A\to e$. For a nonzero associative unital ring $k$ this
specializes to the complete classification
\[
\begin{split}
 \operatorname{Loc}(7F)\text{ for }W_{D(k)}^{\mathrm{lc}}
 &\Longleftrightarrow
  \text{no rational prime is invertible in }k\\
 &\Longleftrightarrow W_{D(k)}^{\mathrm{lc}}=W_\infty.
\end{split}
\tag{B.7.1}
\]
The constant $k$-homology localizer satisfies none of
Loc~(7a,b,c,F), for every nonzero $k$.
For $\Wone$, which need not be a localizer,
Loc~(7c) nevertheless holds for every derivator.
With the enhancement of C.2, Loc~(7a) and the
factorization assertion of Loc~(7F) hold as well;
Loc~(7b) fails for sets, groupoids, every finite
$r$-type theory, every nonzero $D(k)$, every nonzero
$L_E\Sp$, and every nontrivial thin coefficient
theory. These facts do not make $\Wone$ the
weak-equivalence class of a Brown category.
\end{claim}
\begin{proof}
\emph{The exact axioms.}
Chapter XVI, ms.~38--39, states
\[
\begin{array}{ll}
 (7a):&\mathcal F_W\subseteq\operatorname{Fib}_W;\\
 (7b):& f:X\to Y\text{ over }S,\quad
 X\to S,Y\to S\in\mathcal F_W,\quad f\in W\\
 &\hspace{1.5em}\Longrightarrow
    T\times_Sf\in W\text{ for every }T\to S;\\
 (7c):&\mathcal F_W\cap W\subseteq W^{\mathrm u};\\
 (7F):&\text{every functor factors as }p i,
       \quad i\in W,\quad p\in\operatorname{Fib}_W.
\end{array}
\tag{B.7.2}
\]
The word ``cofibrations'' in the manuscript's formulation
of (7F) is a typographical slip: the surrounding definition
is the category-of-fibrant-objects factorization axiom.
The class called \emph{weak $W$-fibrations} in
\citet[Definition~6.1.13]{Cisinski2006} is precisely
$\operatorname{Fib}_W$ here. It is not the manuscript's
weaker class obtained by testing only contractible bases.

\emph{A pullback lemma.}
Suppose that factorizations in (7F) exist. If
$f:X\to Y$ is in $W$ over $S$ and both structural maps
belong to $\operatorname{Fib}_W$, then every pullback
of $f$ over $S$ is in $W$. Indeed, factor $t:T\to S$
as $T\xrightarrow{a}Z\xrightarrow{b}S$, with
$a\in W$ and $b\in\operatorname{Fib}_W$.
The map $f_Z$ is in $W$ by the definition of $b$.
The maps $X_T\to X_Z$ and $Y_T\to Y_Z$ are in $W$
by the definition of the pullbacks of the two structural
maps. Two-out-of-three in the resulting square gives
$f_T\in W$. In particular
\[
 \operatorname{Fib}_W\cap W\subseteq W^{\mathrm u}.
\tag{B.7.3}
\]

\emph{Why (7a) implies (7b).}
Every functor $u:A\to S$ factors as a weak homotopy
equivalence followed by a member of $\mathcal F_{W_\infty}$.
For completeness, factor it first as
$A\to P\xrightarrow{p}S$ in the ordinary Thomason model
structure. Its model fibration $p$ is universally
$W_\infty$-preserving, by right properness. Form the
Grothendieck fibration $q:\Phi(P)\to S$ of the
contravariant diagram $s\mapsto s\backslash P$.
An object is $(s,x,\beta:s\to p(x))$; an arrow to
$(s',x',\beta')$ is a pair $(a:s\to s',h:x\to x')$
with $p(h)\beta=\beta'a$. The inclusion
$j(x)=(p(x),x,\id)$ is a weak homotopy equivalence:
forgetting $s,\beta$ gives a retraction $\epsilon$,
and the arrows $(\beta,\id)$ give
$\id\Rightarrow j\epsilon$.
For $s\to s'$, the transition
$s'\backslash P\to s\backslash P$ is the pullback
by $p$ of $s'\backslash S\to s\backslash S$.
The latter is a weak equivalence between categories
with initial objects, so the transition is a weak
equivalence. This proves the asserted factorization.

Since $\mathcal F_{W_\infty}\subseteq\mathcal F_W$,
(7a) supplies (7F), using these factorizations. Applying
the pullback lemma then gives (7b). Taking $Y=S$ and
$q=\id_S$ gives (7b)$\Rightarrow$(7c).

\emph{Why (7c) forces the minimal localizer.}
The exact hypotheses of
\citet[Theorem~6.4.22]{Cisinski2006} are:
(a) $W\subseteq W_0$; and (b) a $W$-locally constant
functor is a $W$-equivalence if and only if it is
$W$-aspherical. No accessibility hypothesis occurs.
Here ``$W$-locally constant'' means that the transitions
$A/s\to A/s'$ are in $W$.
Given such a functor $u:A\to S$, replace it by the
Grothendieck opfibration
\[
 q:\int_{s\in S}(A/s)\longrightarrow S.
\]
The projection to $A$ is a weak homotopy equivalence,
with section $a\mapsto(u(a),a,\id)$ and the evident
natural comparison; that comparison also gives
$q\in W\Longleftrightarrow u\in W$.
Its fibres are $A/s$ and its transports belong to $W$.
Apply (7c) to its opposite, a Grothendieck fibration.
Fundamental localizers are invariant under opposites.
We conclude that every $A/s$ is $W$-aspherical.
Thus a $W$-locally constant $W$-equivalence is
$W$-aspherical. The converse is Theorem A.
The cited theorem proves $W=W_\infty$.

For $W_\infty$, a Grothendieck fibration with weakly
equivalent transports computes its strict base changes
as homotopy pullbacks, by Quillen's Theorem B and its
base-change criterion
\citep[\S1, Theorem B and its corollary]{Quillen1973};
see also \citet[Proposition~2.7]{Rezk1998Sharp}.
It therefore preserves weak equivalences under every
base change, giving (7a), and hence (7b,c). This completes
their classification. The same result holds for every
nontrivial fundamental localizer, without the explicit
condition $W\subseteq W_0$, by
\citet[Remark~6.4.23]{Cisinski2006}. The only additional
nontrivial case is the gross localizer; the inclusion of
one point into two discrete points refutes (7c) there.

\emph{Why (7F) is exactly properness.}
Let $W_{\mathrm{asph}}$ denote the fundamental localizer
generated by all projections $A\to e$ which belong
to $W$. It is contained in $W$ and has exactly the same
aspherical categories and aspherical functors
\citep[\S6.1.7]{Cisinski2006}.
Suppose (7F). For any $f:X\to Y$ in $W$, factor its
graph as
\[
 X\xrightarrow{i}P\xrightarrow{q}X\times Y,
 \qquad i\in W,\quad q\in\operatorname{Fib}_W.
\]
The projections from a product belong to
$\operatorname{Fib}_W$, since a fundamental localizer
is stable under products. Consequently
$r=\operatorname{pr}_Xq$ and $p=\operatorname{pr}_Yq$
belong to $\operatorname{Fib}_W$. Since $ri=\id_X$
and $pi=f$, both $r,p$ are also in $W$.
By (B.7.3) they are universally in $W$.
For each $y\in Y$, pull $p$ back to $Y/y$. The result
$P/y\to Y/y$ is in $W$, and its target is contractible.
Thus $P/y$ is $W$-aspherical, equivalently
$W_{\mathrm{asph}}$-aspherical. Theorem A for
$W_{\mathrm{asph}}$ gives $p\in W_{\mathrm{asph}}$.
The same argument gives $r\in W_{\mathrm{asph}}$.
Then $ri=\id_X$ gives $i\in W_{\mathrm{asph}}$, and
$f=pi\in W_{\mathrm{asph}}$. We have proved
\[
 \operatorname{Loc}(7F)\Longrightarrow W=W_{\mathrm{asph}},
\tag{B.7.4}
\]
without an accessibility assumption.

If $W$ is accessible, Proposition~6.1.8 and
Theorem~6.1.11 of \citet{Cisinski2006} identify this
equality with properness, or equivalently generation
by a set of aspherical projections. Conversely, for
proper $W$ the Thomason model structure is right proper
\citep[Theorem~5.2.15]{Cisinski2006}. Its model
fibrations belong to $\operatorname{Fib}_W$
\citep[Example~6.1.14]{Cisinski2006}, and its
factorization axiom supplies (7F).
This proves the claimed equivalence.
Here properness is right properness of the associated
simplicial model structure; the cited theorem gives the
equivalent Thomason-model condition. Literal Loc~(9),
accessibility of the full subcategory of arrows, must not
be confused with the additional assertion
$W=W_{\mathrm{asph}}$.

\emph{An obstruction to properness.}
We use the precise statement of
\citet[Lemma~9.4.5]{Cisinski2006}: if $W\subseteq W_n$
and $W$ is proper, every simplicial $W$-equivalence
whose target is $W$-fibrant is a \emph{strong}
$n$-equivalence. By Scholie~9.2.18 there, this includes
surjectivity on $\pi_{n+1}$, in addition to isomorphisms
through degree $n$. It is enough below to exhibit a
$W$-equivalence between connected spaces, with
$W$-local target and nonsurjective $\pi_{n+1}$.
A Kan representative of a $W$-local space is
$W$-fibrant: restriction against a $W$-equivalence
which is a cofibration is a weak equivalence and a
Kan fibration of mapping spaces, hence a trivial
fibration. This verifies the target hypothesis in
each application, rather than omitting it.

\emph{Characteristic zero rings.}
Let $S$ be the set of rational primes invertible in
$k$, and set $R=\Z[S^{-1}]$. By the arithmetic test
of C.9, $k$-homology equivalences are exactly
$R$-homology equivalences: both are detected by
$\Q$ and by $\F_p$ for $p\notin S$.
If $S$ is empty, C.9 gives
$W_{D(k)}^{\mathrm{lc}}=W_\infty$, which is proper.
If $S$ is nonempty, consider
\[
 j:K(\Z,2)\longrightarrow K(R,2).
\tag{B.7.5}
\]
Its homotopy fibre is $K(R/\Z,1)$. The group $R/\Z$
is the filtered union of finite groups whose orders
are invertible in $R$. Averaging makes their positive
$R$-homology zero, and the bar construction commutes
with filtered unions. Thus this fibre is $R$-acyclic.
The Serre homology spectral sequence proves that $j$
is an $R$-homology equivalence. Both its spaces are
simply connected, so C.7 makes it a
$W_{D(k)}^{\mathrm{lc}}$-equivalence.

The target is local for $R$-homology equivalences.
Indeed, the homotopy groups of
$\Map(X,K(R,2))$ are $H^{2-i}(X;R)$ for $0\leq i\leq2$;
these are invariant under $R$-homology equivalences,
by applying derived $R$-linear Hom into $R$ to
the singular-chain equivalence. It is consequently
local for the smaller class $W_{D(k)}^{\mathrm{lc}}$.
But $\pi_2(j):\Z\to R$ is not onto. Since C.7
also gives $W_{D(k)}^{\mathrm{lc}}\subseteq W_1$,
the preceding obstruction with $n=1$ proves that
this localizer is not proper.

\emph{Positive characteristic rings.}
Let $d=\operatorname{char}(k)>0$, and let $P$ be
the nonempty finite set of primes dividing $d$.
The nonunits among rational primes are precisely
the members of $P$. C.9 identifies $k$-homology
equivalences with simultaneous $\F_p$-homology
equivalences for $p\in P$. Put
\[
 M=\prod_{p\in P}\Z_p,
 \qquad j:K(\Z,2)\longrightarrow K(M,2),
\tag{B.7.6}
\]
where $\Z\to M$ is diagonal and $\Z_p$ denotes
the additive group of $p$-adic integers, regarded here as a discrete group.
For each $p\in P$, multiplication by $p$ is an
automorphism of $A=M/\Z$. For surjectivity, given
$m\in M$, choose an integer $z$ congruent to its
$p$-component modulo $p$. Then $(m-z)/p\in M$:
division by $p$ is allowed in the other components.
For injectivity, if $pm$ is the diagonal integer
$z$, integrality in the $p$-component implies
$p\mid z$, and then $m$ is itself diagonal.

For clarity, an abelian group uniquely divisible
by $p$ is $\F_p$-acyclic as a group. Such a group
has no $p$-torsion. Each finitely generated subgroup
$H$ is therefore a sum of a free abelian group and
a finite group of order prime to $p$, so
$H_*(BH;\F_p)$ is the exterior algebra on its first
homology. Enlarge $H$ to the finitely generated
subgroup $p^{-1}H$. Its inclusion induces zero on
$H_1(-;\F_p)$, since under the isomorphism
$p^{-1}H\simeq H$ it is multiplication by $p$.
It therefore induces zero in every positive
homology degree. Filtered colimits of the bar
complexes prove the assertion. Hence $K(A,1)$,
the homotopy fibre of (B.7.6), is $\F_p$-acyclic
for every $p\in P$. The same Serre argument proves
that $j$ is in $W_{D(k)}^{\mathrm{lc}}$.

The target in (B.7.6) is local for $k$-homology
equivalences. To verify this rather than assume
it, an $\F_p$-homology equivalence induces all
$\F_p$-cohomology isomorphisms by vector-space
duality. The coefficient short exact sequences
\[
 0\longrightarrow\Z/p^{r-1}
 \longrightarrow\Z/p^r\longrightarrow\F_p
 \longrightarrow0
\]
and their long exact cohomology sequences imply
the same assertion for $\Z/p^r$. Thus
$K(\Z/p^r,2)$ is $\F_p$-homology local.
Local spaces are closed under homotopy limits,
and the surjective inverse system of these
coefficient groups gives
\[
 \holim_rK(\Z/p^r,2)\simeq K(\Z_p,2);
\]
the possible derived-limit group is zero by
surjectivity. Taking the finite product over $P$
proves the required locality for the simultaneous
homology class, and hence for
$W_{D(k)}^{\mathrm{lc}}$. The map
$\pi_2(j):\Z\to M$ is not surjective, since $M$
is uncountable. The $n=1$ obstruction again
proves nonproperness. Together with B.4 and
(B.7.4), this proves (B.7.1) for every nonzero ring.

\emph{The constant homology class.}
For every nonzero $k$, it contains $BG\to e$ for
the nontrivial acyclic group used in B.6, so it
strictly contains $W_\infty$. It preserves
components, and the already proved classification
therefore refutes (7a,b,c). It is also never proper.
In characteristic zero with $S$ empty it equals
the integral homology localizer, whose
nonproperness is
\citet[Scholie~9.4.7]{Cisinski2006}.
When $S$ is nonempty, use
$K(\Z,1)\to K(R,1)$ and the obstruction with
$n=0$. Its target is $R$-homology local by the
same cohomology argument. Its $R$-homology is an
isomorphism: $R$ is the filtered union of its
cyclic subgroups, the bar homology commutes with
that union, and the only positive group homology
of $\Z$ is in degree one. The map on $\pi_1$ is
not surjective.

In positive characteristic use
$K(\Z,1)\to K(M,1)$. The preceding locality proof
works in degree one as well. The torsion-free
abelian group $M$ is a filtered union of finitely
generated free abelian groups, whence
\[
 H_*(BM;\F_p)\simeq
 \bigwedge_{\F_p}^*(M/pM)
 \simeq\bigwedge_{\F_p}^*\F_p
 \quad(p\in P).
\]
The diagonal $\Z\to M$ induces the displayed
isomorphism in degree one, and hence is a
$k$-homology equivalence. Again it is not onto
on $\pi_1$. Lemma~9.4.5 with $n=0$ proves
nonproperness. Accessibility of these homology
localizers is B.4, or
\citet[Theorem~9.4.2 and Corollary~9.4.3]{Cisinski2006}.
Consequently (7F) fails in every nonzero-ring case.

\emph{The intermediate class $\Wone$.}
The fundamental-localizer implications above must
not be applied to this class, because it need not
satisfy two-out-of-three. We instead check each
assertion directly. For variance, put
$\D^{\op}(A)=\D(A^{\op})^{\op}$. Then
$f\in W^1_{\D}$ if and only if
$f^{\op}\in W^1_{\D^{\op}}$.
It therefore suffices to prove the assertion for
Grothendieck opfibrations and apply that proof
to $\D^{\op}$ when the manuscript uses
Grothendieck fibrations. This dualizes both the
shape and coefficient derivator; it does not
assume shape-opposite invariance for a fixed
bare derivator. For opfibrations, right Kan
extensions in the presheaf convention have
the strict-fibre formula of B.3.
If $p:E\to S$ has transports in $\Wone$, then
$p_*$ sends locally constant objects to locally
constant objects, and does so after every base
change. Indeed, each transport is in the
cohomological-equivalence class by A.3; the
comparison of the two fibre cohomology objects
is an isomorphism for every locally constant
coefficient. Pointwise Kan extension gives
exactly this comparison on the base arrow.
Consequently $p^*$ and $p_*$ restrict to an
adjunction on locally constant categories.

If also $p\in\Wone$, the unit of this adjunction
is an isomorphism. Apply it to the constant
object of value $M\in\D(e)$ and evaluate at
$s\in S$. The strict-fibre formula gives
\[
 M\xrightarrow{\sim}
 (p_{E_s})_*p_{E_s}^*M.
\tag{B.7.8}
\]
Thus $E_s\to e$ is a constant-coefficient
equivalence, equivalently a member of $\Wone$
for this point target, by A.3. Every base change
of $p$ is again an opfibration with the same
fibres. The fully faithful case of B.3 proves
that each such base change belongs to $\Wone$.
This is literal Loc~(7c), without any two-out-of-three
argument.

Under the enhancement, C.2 identifies
$\Wone$ with the cohomological class. The
all-local-coefficients direct-image property
just proved and its Beck--Chevalley formula
make $p$ universally preserve that class, by
the adjunction argument of D.2. Explicitly,
for $f:T\to S$ in the cohomological class and
$F\in\LC(E)$, cohomology on $E$ equals
cohomology on $S$ with coefficient $p_*F$.
This coefficient is locally constant, so $f$
preserves it; base change identifies the result
with the comparison on $T\times_SE$.
This proves Loc~(7a) for $\Wone$ in the
enhanced case. For the factorization assertion,
factor any functor in the ordinary Thomason
model structure. Its first map is in
$W_\infty\subseteq\Wone$. Its second map
is an ordinary sharp fibration, whose strict
base changes compute homotopy pullbacks. C.2
makes $\Wone$ stable under homotopy pullback,
so this second map is in
$\operatorname{Fib}_{\Wone}$. This proves
the asserted factorization, without claiming
the Brown axioms that require two-out-of-three.

Finally choose a pointed category $Z$ with
$Z\to e\in\Wone$ and $e\to Z\notin\Wone$.
Let $R=Z\star\{t\}$ and $p:R\to[1]$ be
the Grothendieck opfibration of $Z\to e$.
The chosen point gives a section
$s:[1]\to R$. Both structural maps $p$ and
$\id_{[1]}$ have transports in $\Wone$.
The section is a weak homotopy equivalence
because both total categories have terminal
objects, but its fibre at $0$ is the excluded
map $e\to Z$. Taking opposites gives the
literal Grothendieck-fibration version. This
refutes (7b), not (7c): the structural map $p$
itself is universally in $\Wone$ by (B.7.8).
For $r$-types take a category model of
$K(A,r+1)$ with $A\ne0$; for groupoids use
$r=1$. For nonzero stable coefficients take
$Z=BG$ for the nontrivial integrally acyclic
group of B.10: its suspension is contractible,
so constants satisfy the fibre criterion C.2,
whereas $e\to BG$ fails full faithfulness by
the regular-action argument in B.6. For a
nontrivial thin coefficient category take
$Z=e\amalg e$ and use (B.6.6).
All these choices meet the two stated tests.
For spaces $\Wone=W_\infty$ and (7b) holds.
\end{proof}

Here is the resulting comparison for the computed theories.
The zero coefficient theory is understood separately: all
its arrows are equivalences and all listed axioms hold.
\begin{center}
\begin{tabular}{p{0.39\linewidth}ccc}
\hline
Class & Loc (6b,c) & Loc (7a,b,c) & Loc (7F)\\
\hline
$W^{\mathrm{lc}}$ for spaces, $D(\Z)$, spectra
 & no & yes & yes\\
$W^{\mathrm{lc}}$ for sets, groupoids, finite $r$-types
 & no & no & yes\\
$W^{\mathrm{lc}}_{D(k)}$, $k\ne0$, all primes nonunits
 & no & yes & yes\\
$W^{\mathrm{lc}}_{D(k)}$, $k\ne0$, some prime a unit
 & no & no & no\\
Constant $k$-homology, $k\ne0$
 & yes & no & no\\
\hline
\end{tabular}
\end{center}
For sets, groupoids and finite $r$-types the last column
uses C.4 and the properness of every finite truncation
localizer $W_n$, explicitly established in
\citet[\S9.2.1]{Cisinski2006}. Thus (7F) does not imply
(7a,b,c). For a general spectrum $E$, B.4 makes
$W^{\mathrm{lc}}_{L_E\Sp}$ accessible, so the exact
remaining test for (7F) is
\[
 W^{\mathrm{lc}}_{L_E\Sp}
   =(W^{\mathrm{lc}}_{L_E\Sp})_{\mathrm{asph}},
\tag{B.7.7}
\]
equivalently properness. No assertion that (B.7.7)
holds for every $E$ is made; its evaluation beyond
the coefficient theories computed here remains
unresolved in this work. For an arbitrary derivator
without an accessibility theorem, (7F) still implies
aspherical generation by (B.7.4); the converse is
not claimed without that hypothesis.

For an arbitrary unenhanced derivator, two intermediate-class
questions remain unresolved here: whether every Grothendieck
fibration with transports in $\Wone$ belongs to
$\operatorname{Fib}_{\Wone}$ (literal Loc~(7a)), and
whether every functor factors as a member of $\Wone$
followed by a member of $\operatorname{Fib}_{\Wone}$
(the factorization assertion of Loc~(7F)). The first
argument above proves universal preservation of the
cohomological class; identifying its output with
$\Wone$ used the enhancement in C.2. The second
argument used the enhanced homotopy-fibre criterion
to prove homotopy-pullback stability of $\Wone$.
Neither of these two extra steps is supplied just
by the arbitrary-derivator relative Theorem A of B.3.
This precisely limits the two positive enhanced
assertions; it does not leave (7c) unresolved.

Finally, the existence of the model-fibration class
in D.6 is compatible with all the negative entries.
An accessible fundamental localizer always has the
Thomason model structure. Its model fibrations compose,
are stable under strict base change, and detect fibre
equivalences between fibrant objects of a slice.
When the localizer is not proper, they need not belong
to the manuscript's $\operatorname{Fib}_W$.
Replacing that distinction by an identification would
incorrectly turn D.6 into Loc~(7F).
